\section*{第 \ref{ch:g8:voltage-voltmeter} 章——电压与电压表}

\begin{solution}{exo:g8:voltage-voltmeter:1}
两点之间的\omterm{def:g8:voltage-voltmeter:voltage}{电压}，就是电路中的“推动力”，
好比瀑布上下两端的落差推动水车。
\omterm{def:g8:voltage-voltmeter:voltage}{电压}是一个差值，因此需要两个点；
其\omterm{def:g6:measurement-in-science:measuring}{单位}是\omterm{def:g8:voltage-voltmeter:voltage}{伏特}，\unit{V}。
\end{solution}

\begin{solution}{exo:g8:voltage-voltmeter:2}
\qty{1.5}{V}; \qty{4.5}{V}; \qty{9}{V}; \qty{12}{V};
\qty{230}{V}.
\end{solution}

\begin{solution}{exo:g8:voltage-voltmeter:3}
完全没有电流——没有闭合回路。\omterm{def:g8:voltage-voltmeter:voltmeter}{电压表}读数为全部的
\qty{1.5}{V}：有\omterm{def:g8:voltage-voltmeter:voltage}{电压}而无电流——推动力可以等待；
没有它，电流便无从开始。
\end{solution}

\begin{solution}{exo:g8:voltage-voltmeter:4}
\omterm{def:g8:intensity-ammeter:ammeter}{电流表}：\omterm{def:g4:series-circuits:series}{串联}接入，相当于为它断开道路，阻碍几乎为零。
\omterm{def:g8:voltage-voltmeter:voltmeter}{电压表}：\omterm{def:g5:series-parallel-circuits:parallel}{并联}接入，不断开道路，阻碍几乎为无穷大。
代价高昂的错误：把\omterm{def:g8:intensity-ammeter:ammeter}{电流表}跨接在两端——\omterm{def:g5:series-parallel-circuits:parallel}{并联}——
就成了带表盘的\omterm{def:g7:short-circuits-safety:device}{短路}。
\end{solution}

\begin{solution}{exo:g8:voltage-voltmeter:5}
器件设计所对应的推动力。在\qty{1.5}{V}下：微弱的闪烁。
在\qty{6}{V}下：达到设计的亮度与寿命。
在\qty{12}{V}下：瞬间强光，随即灯丝烧断。
\end{solution}

\begin{solution}{exo:g8:voltage-voltmeter:6}
\omterm{def:g3:first-electric-circuit:battery}{电池}：\qty{4.5}{V}；\omterm{def:g3:first-electric-circuit:bulb}{灯泡}：\qty{4.5}{V}；导线：约为\qty{0}{V}。
推动力存在于用电器两端——
导线则是一条平坦的道路。
\end{solution}

\begin{solution}{exo:g8:voltage-voltmeter:7}
熄灭的\omterm{def:g3:first-electric-circuit:bulb}{灯泡}：\qty{0}{V}；断开的\omterm{ex:g3:first-electric-circuit:switch}{开关}：全部\qty{4.5}{V}，
等待于断口两端。\omterm{ex:g3:first-electric-circuit:switch}{开关}闭合的瞬间，二者互换角色：
闭合的\omterm{ex:g3:first-electric-circuit:switch}{开关}两端降至约为零，发光的\omterm{def:g3:first-electric-circuit:bulb}{灯泡}承担全部
\qty{4.5}{V}。
\end{solution}

\begin{solution}{exo:g8:voltage-voltmeter:8}
它拒绝让电流通过：阻碍几乎完全，自身几乎不分流，
因此\omterm{def:g3:first-electric-circuit:bulb}{灯泡}中的电流与发光照常进行，
仿佛从未被观察到。
\end{solution}

\begin{solution}{exo:g8:voltage-voltmeter:9}
三节\omterm{def:g3:first-electric-circuit:battery}{电池}首尾相接：$1.5 + 1.5 + 1.5 = \qty{4.5}{V}$。一节反接，
与另外两节相互抵消：$1.5 + 1.5 - 1.5 = \qty{1.5}{V}$——
这正是多年前手电筒的谜题，如今用\omterm{def:g8:voltage-voltmeter:voltage}{伏特}来解释。
\end{solution}

\begin{solution}{exo:g8:voltage-voltmeter:10}
新\omterm{def:g3:first-electric-circuit:battery}{电池}：略高于\omterm{def:g8:voltage-voltmeter:nominal}{额定电压}——亮度饱满。
随着使用时间增加，\omterm{def:g3:first-electric-circuit:battery}{电池}\omterm{def:g8:voltage-voltmeter:voltage}{电压}降至\qty{11.2}{V}，
\omterm{def:g3:first-electric-circuit:bulb}{灯泡}亮度略低于设计的亮度。
没有损坏：\omterm{def:g8:voltage-voltmeter:voltage}{电压}只会\emph{低于}额定值；
这一约定惩罚的是超出，而不是不足。
\end{solution}

\begin{solution}{exo:g8:voltage-voltmeter:11}
\omterm{def:g8:voltage-voltmeter:voltmeter}{电压表}串在道路中，几乎完全阻碍电流：
电流几乎停止，\omterm{def:g3:first-electric-circuit:bulb}{灯泡}不亮。
\omterm{def:g8:voltage-voltmeter:voltmeter}{电压表}跨接在回路唯一的推动力两端，
读数接近\omterm{def:g3:first-electric-circuit:battery}{电池}的全部\omterm{def:g8:voltage-voltmeter:voltage}{电压}——对于接错的电路，这是正确的读数：
很有启发意义，却不会起火。
\end{solution}

\begin{solution}{exo:g8:voltage-voltmeter:12}
\omterm{def:g5:series-parallel-circuits:parallel}{并联}时，每只\omterm{def:g3:first-electric-circuit:bulb}{灯泡}直接跨接在\omterm{def:g3:first-electric-circuit:battery}{电池}两端：每只均为\qty{4.5}{V}，
满额定推动力、满亮度，正如旧的\omterm{def:g5:series-parallel-circuits:parallel}{并联}规则所预言的那样。
\omterm{def:g4:series-circuits:series}{串联}时，两只\omterm{def:g3:first-electric-circuit:bulb}{灯泡}\emph{分享}\qty{4.5}{V}——每只约为\qty{2.25}{V}，
两只\omterm{def:g3:first-electric-circuit:bulb}{灯泡}都分得不足的\omterm{def:g8:voltage-voltmeter:voltage}{电压}，发光暗淡，
正如旧的\omterm{def:g4:series-circuits:series}{串联}规则所观察到的那样。
隔一章之后，这种分配将被写成\omterm{def:g8:voltage-voltmeter:voltage}{电压}的规律。
\end{solution}

\begin{solution}{pb:g8:voltage-voltmeter:1}
\textbf{1.} 将\omterm{def:g8:voltage-voltmeter:voltmeter}{电压表}的两根表笔分别接在\omterm{def:g3:first-electric-circuit:battery}{电池}的两个\omterm{def:g3:first-electric-circuit:battery}{接线柱}上——
跨接，即\omterm{def:g5:series-parallel-circuits:parallel}{并联}；\omterm{def:g8:voltage-voltmeter:voltmeter}{电压表}不切断电路，也不需要单独的电路：
它\omterm{def:g6:measurement-in-science:measuring}{测量}的是等待中的推动力。
\textbf{2.} \qty{1.58}{V}：新\omterm{def:g3:first-electric-circuit:battery}{电池}。\qty{1.32}{V}：仍可用。
\qty{0.9}{V}：电量不足。\qty{0.1}{V}：已耗尽——从抽屉丢进垃圾桶。
\textbf{3.} 每节约为 $8.9 \div 6 \approx \qty{1.48}{V}$——
六节完好的\omterm{def:g3:first-electric-circuit:battery}{电池}。
\textbf{4.} 它的化学物质仍能产生微弱的推动力，但已无法
维持有用的电流：\omterm{def:g8:voltage-voltmeter:voltage}{电压}只是一个承诺，
在负载下，\omterm{def:g3:first-electric-circuit:battery}{电池}已没有力量兑现。
\textbf{5.} 两节首尾相接：应得 $2 \times 1.5 =
\qty{3}{V}$；剩下的两节只提供 $2 \times 1.32 =
\qty{2.64}{V}$——相机或许还能工作，但会有些吃力。
\textbf{6.} 四节\omterm{def:g3:first-electric-circuit:battery}{电池}的总\omterm{def:g8:voltage-voltmeter:voltage}{电压}须达到\qty{5}{V}或更高：四节新\omterm{def:g3:first-electric-circuit:battery}{电池}
各为 $1.58$，总计\qty{6.3}{V}——足够；混入 $1.32$ 的\omterm{def:g3:first-electric-circuit:battery}{电池}时，
只要总\omterm{def:g8:voltage-voltmeter:voltage}{电压}至少为\qty{5}{V}即可（例如两节新的与两节仍可用的：
$1.58 \times 2 + 1.32 \times 2 =
\qty{5.8}{V}$）。电量不足和已耗尽的\omterm{def:g3:first-electric-circuit:battery}{电池}不宜使用。
\textbf{7.} 三节\omterm{def:g3:first-electric-circuit:battery}{电池}：\qty{4.5}{V}接到\qty{3}{V}的\omterm{def:g3:first-electric-circuit:bulb}{灯泡}上——
超出额定值一半：至多是一个过亮而光彩夺目的夜晚，
最可能的结果是\omterm{def:g3:first-electric-circuit:bulb}{灯泡}很快烧毁。
这一约定惩罚的是超出。
\textbf{8.} 缓慢的损伤，温和却真实：\qty{9}{V}的\omterm{def:g8:voltage-voltmeter:voltage}{电压}加在
\qty{8}{V}的器件上，就是持续的过压——它仍能工作，但会发热，
并且迅速老化。额定值就是额定值；
“差不多”正是灯丝早夭的原因。
\textbf{9.} \omterm{def:g3:first-electric-circuit:battery}{电池}的推动力完整地等待于唯一的真实缺口——断开的\omterm{ex:g3:first-electric-circuit:switch}{开关}两端；
空载的\omterm{def:g3:first-electric-circuit:bulb}{灯泡}不通过电流，因此其两端不显示\omterm{def:g8:voltage-voltmeter:voltage}{电压}。
\omterm{ex:g3:first-electric-circuit:switch}{开关}闭合后：\omterm{ex:g3:first-electric-circuit:switch}{开关}两端接近\qty{0}{V}，\omterm{def:g3:first-electric-circuit:bulb}{灯泡}两端接近
\qty{3.1}{V}——推动力从缺口
转移到了用电器上。
\textbf{10.} 这根良好的导线是一条平坦的道路：它输送电流，
却不消耗推动力——所有\omterm{def:g8:voltage-voltmeter:voltage}{电压}都留给了
用电器。
\textbf{11.} “插座\omterm{def:g8:voltage-voltmeter:voltage}{电压}为\qty{230}{V}，是这个抽屉里最强\omterm{def:g3:first-electric-circuit:battery}{电池}的
五十倍；推动力放大五十倍，同样潮湿的双手中
就会流过五十倍的电流，心脏会在百分之几\omterm{def:g8:intensity-ammeter:intensity}{安培}时
停止跳动。事故的结局，
就在这个抽屉里。”
\textbf{12.} 例如：“\omterm{def:g8:voltage-voltmeter:voltmeter}{电压表}总是\omterm{def:g5:series-parallel-circuits:parallel}{并联}在电路两端，
绝不串在电路中。标签上的\omterm{def:g8:voltage-voltmeter:nominal}{额定电压}既承诺了设计的性能，
也警示超过该值将导致烧毁。
而在任何\omterm{def:g3:first-electric-circuit:circuit}{断开的电路}中，\omterm{def:g8:voltage-voltmeter:voltage}{电压}都等待于缺口之处——
跨在断开的\omterm{ex:g3:first-electric-circuit:switch}{开关}两端，随时待命。”
\end{solution}
